\section{Post-Training 数据集：质量的三个维度}

\subsection{好数据的定义}

讲者认为数据质量是整个 Post-Training 流程中最核心的问题，可以从三个维度来衡量：

\begin{importantbox}{数据质量的三维框架}
\begin{enumerate}
    \item \textbf{Accuracy（准确性）}：回答是否正确。对于简单问题较易验证；对于代码可通过 unit test 验证，对于数学可通过 solver 验证。
    \item \textbf{Diversity（多样性）}：样本之间是否覆盖了用户可能的各种交互场景。多样性不足会导致模型在某些领域能力严重缺失。
    \item \textbf{Complexity（复杂度）}：样本是否足够有挑战性，包含多步推理、chain-of-thought 等。
\end{enumerate}
\end{importantbox}

\subsection{多样性的重要性}

通过对不同数据集的 embedding 进行 2D 可视化，可以直观看到差异：

\begin{itemize}
    \item \textbf{真实用户对话}（如 ShareGPT、WildChat）：embedding 分布覆盖极广，因为用户提问涉及方方面面
    \item \textbf{数学数据集}：embedding 高度聚集在单一区域
    \item \textbf{代码数据集}：同样集中在有限区域
\end{itemize}

因此，若目标是构建通用 LLM，仅有数学和代码数据远远不够，必须引入多样化的对话数据来扩展覆盖范围。

\begin{warningbox}{合成数据的多样性陷阱}
合成数据（synthetic data）的一大风险在于多样性坍缩。少量合成数据可以有效扩充多样性，但\textbf{过多}合成数据反而会导致数据分布收窄、多样性急剧下降，最终降低整体数据质量。
\end{warningbox}

\subsection{复杂度提升：Auto-Evol}

低复杂度样本如"埃菲尔铁塔多高？"只需几个词即可回答。可以使用 \textbf{Auto-Evol} 等框架，让 LLM 对 prompt 进行改写以提升复杂度——例如要求多角度分析、比较不同时期数据、解释工程原理等。改写后的问题需要更长、更详细的回答，自然提升了数据的复杂度。

\subsection{数据格式}

Post-Training 使用两种主要数据格式：

\begin{table}[H]
\centering
\caption{Post-Training 数据格式}
\begin{tabular}{lll}
\toprule
\textbf{格式} & \textbf{组成} & \textbf{用途} \\
\midrule
Instruction Data & system prompt + instruction + output & SFT \\
Preference Data & instruction + chosen answer + rejected answer & DPO/PPO \\
\bottomrule
\end{tabular}
\end{table}

\subsection{本章小结}

数据质量由准确性、多样性和复杂度三个维度共同决定。在构建数据集时，需特别关注合成数据的多样性风险，并可使用 Auto-Evol 等工具提升样本复杂度。

\newpage
